\section{Early Activation on the Seeded Contrast}
\label{sec:phil-early-parent}

The two preceding sections activate the channel of
\cref{prop:rg-formation-endogenous-cycle}(b) after the late restored
coherence of the joint slow schedule. This section returns to the early
contrast of \cref{prop:phil-seeded-contrast-formation} at every mobility
\(\epsilon_{\mathrm L}\in[0,1]\), a range that contains the original coupled
mobility \(\epsilon_{\mathrm L}=\mathrm T^{1/8}\), and conditions the
recognition law on the seeded sub-event of that contrast.
\Cref{prop:phil-early-seeded-parent} bounds the discarded conditional
information at a deterministic activation time from above and from below,
\cref{prop:phil-early-first-completion} treats activation at the first
qualifying completion time of the first block, and
\cref{prop:phil-early-parent-dynamics} couples the retained path of the
own-law parent with the conditioned native path.
\Cref{rem:phil-early-stationary-mismatch} records why the stationary
comparison of \cref{prop:phil-actual-parent-dynamics} is not used.

\definitionheading{Seeded early prefix}{def:phil-early-seeded-prefix}
Take the model of \cref{conj:phil-contrast-formation} at the temperature
\(\mathrm T=\epsilon^8\) of \cref{eq:rg-formation-cycle-regime}, with
\(0<\epsilon\le1/100\), a learning mobility \(\epsilon_{\mathrm L}\ge0\), and
the initial law \(\mathcal N(\mathsf m_0,I_{12})\) times the uniform law on
assignments; \(\epsilon_{\mathrm L}=0\) is the frozen arm and
\(\epsilon_{\mathrm L}=\epsilon\) is the regime of
\cref{prop:phil-seeded-contrast-formation}. Write \(\pi_{ij}(t)\) for the
receiver rows, \(p_i=\pi_{i,i+1}\), \(K_S=\nabla^2U_S\) for the drift matrix
of sheet \(S\), written \(\mathsf A_S\) in
\cref{prop:phil-actual-parent-dynamics},
\(\mu_{\mathrm T,p}\propto\pi_S(p)e^{-U_S/\mathrm T}\) for the
restricted optimizer at rows \(p\), and
\[
 \Lambda_{\mathrm L}(t)=2e^{\epsilon_{\mathrm L}t}-1,\qquad
 \delta_\epsilon=\frac{\epsilon^2}{8\sqrt3},\qquad
 \theta_\circ=\tfrac12\log(1/\epsilon)+1 .
\]
The kinetics are driven by mutually independent inputs: the initial state
\((\mathbf X_0,S_0)\), a standard Brownian motion \(W\) on \(\R^{12}\), a
rate-one clock, and uniform variables on \([0,1]\), one for each ring and
receiver. At a ring at time \(s\), receiver \(i\) takes source \(i+1\)
exactly when its variable is at most \(\beta_{i,i+1}(\mathbf X_{s-};\pi(s))\),
where
\[
 \beta_{ij}(\mathbf x;\pi)=
 \frac{\pi_{ij}e^{-\mathcal D_{e(i,j)}(\mathbf x)/\mathrm T}}
      {\sum_{j'=i\pm1}\pi_{ij'}e^{-\mathcal D_{e(i,j')}(\mathbf x)/\mathrm T}},
\]
and \(\mathcal F_t\) is generated by the initial state and the inputs on
\([0,t]\). For a perfect matching \(M\in\{S^{(0)},S^{(1)}\}\) with edges
\(e,f\), let \(F_M\) have the orthonormal columns
\((n^+_e-n^+_f)/\sqrt2\), \(n^-_e\), \(n^-_f\) of
\cref{eq:rg-formation-cycle-normals}, and put \(Z_M=\operatorname{im}F_M\),
\(W_M=Z_M^\perp\),
\(K_Z=\operatorname{diag}(8+\mathrm T,4+5\mathrm T,4+5\mathrm T)\), and
\(\nu_{\mathrm T,Z}=\mathcal N(0,\mathrm TK_Z^{-1})\). The hidden coordinates
are \(z=F_M^\top\mathbf x=(z_1,z_2,z_3)\), so that
\(z_2=n^-_e\cdot\mathbf x\) and \(z_3=n^-_f\cdot\mathbf x\), and the retained
common coordinate is \(z_+=u_M\cdot\mathbf x\) with
\(u_M=(n^+_e+n^+_f)/\sqrt2\). The held process, its noise, and the center
path are
\[
 \widetilde{\mathbf X}^M_t=e^{-K_Mt}\mathbf X_0+\mathbf Z^M_t,\qquad
 \mathbf Z^M_t=\sqrt{2\mathrm T}\int_0^te^{-K_M(t-s)}dW_s,\qquad
 \bar{\mathbf x}^{\sigma,M}(t)=\sigma e^{-K_Mt}\mathbf x^\circ,
\]
for \(\sigma\in\{1,-1\}\). The seeds, remainders, and seeded prefix are
\[
\begin{aligned}
 \mathscr B_{\sigma,M}&=\{|\mathbf X_0-\sigma\mathbf x^\circ|<\varrho_\circ\}
   \cap\{S_0=M\}\cap\{\text{no ring in }[0,1/4]\},\\
 \mathscr A^M_t&=\Bigl\{\sup_{s\le t}|\mathbf Z^M_s|\le\delta_\epsilon\Bigr\}
   \cap\{S_s=M\text{ for all }s\le t\},\\
 \mathscr E^\circ_\theta&=\bigcup_{\sigma,M}
   \bigl(\mathscr B_{\sigma,M}\cap\mathscr A^M_\theta\bigr),
\end{aligned}
\]
and \(\mathscr B^\circ=\bigcup_{\sigma,M}\mathscr B_{\sigma,M}\). At a
deterministic activation time \(\theta\) the conditioned law is
\(q^\circ_\theta=\operatorname{Law}(\mathbf X_\theta,S_\theta\mid
\mathscr E^\circ_\theta)\), the channel \(C\) is that of
\cref{prop:rg-formation-endogenous-cycle}(b), and
\(\mu_{\mathrm{act}}=C_\#q^\circ_\theta\). The functionals
\(\mathcal H_{\mathrm{opt}}\) and \(\mathcal H_{\mathrm{phys}}\) are those of
\cref{prop:phil-actual-conditioned-parent}, and \(\mathcal S_2\) is the binary
entropy of \cref{eq:rg-formation-cycle-constants}.

\lemmaheading{Discarded information under conditioning and mixing}{lem:phil-early-discarded-information}
Let \((Y,Z)\) be a random element of a product of standard Borel spaces, with
law \(Q\), let \(\kappa\) be a Markov kernel from the space of \(Y\) to that
of \(Z\), and put
\[
 \mathcal H_\kappa(Q)=\E\KL\bigl(\operatorname{Law}(Z\mid Y)\Vert\kappa_Y\bigr)
 \in[0,\infty].
\]
(i) Let \(B\) be a random variable with finitely many values \(b\),
probabilities \(w_b>0\), and conditional laws \(Q_b\) of \((Y,Z)\) given
\(B=b\). Then
\[
 \sum_bw_b\mathcal H_\kappa(Q_b)=\mathcal H_\kappa(Q)+\mathbb I(B;Z\mid Y),
 \qquad
 0\le\mathbb I(B;Z\mid Y)\le-\sum_bw_b\log w_b .
\]
In particular \(\mathcal H_\kappa\) is convex. If \(A\) is an event with
\(\Pr(A^c)\le\bar r<1/2\) and \(Q_A\) is the conditional law of \((Y,Z)\)
given \(A\), then
\(\mathcal H_\kappa(Q_A)\le[\mathcal H_\kappa(Q)+\mathcal S_2(\bar r)]/(1-\bar r)\),
and for every event \(A\) of positive probability,
\(\mathcal H_\kappa(Q)\ge\Pr(A)\mathcal H_\kappa(Q_A)-\log2\).

(ii) If \(\Upsilon\) is a random element of a standard Borel space defined
jointly with \((Y,Z)\), then
\(\mathcal H_\kappa(Q)\le\E\KL(\operatorname{Law}(Z\mid Y,\Upsilon)\Vert\kappa_Y)\).

(iii) If \(\mathrm g\ge0\) is measurable and
\(\int e^{\mathrm g(y,z)}\kappa_y(dz)\le\mathrm c_1\) for every \(y\), then
\(\mathcal H_\kappa(Q)\ge\E\mathrm g(Y,Z)-\log\mathrm c_1\).

For a law \(q\) of \((\mathbf X,S)\) in the four-cycle,
\(\mathcal H_{\mathrm{opt}}(q)=\mathcal H_\kappa\) of the law of
\((C(\mathbf X,S),Z)\), where \(Z=F_S^\top\mathbf X\) and
\(\kappa=\nu_{\mathrm T,Z}\) on the matching strata and \(Z=0\),
\(\kappa=\delta_0\) elsewhere; \(\mathcal H_{\mathrm{phys}}(q)\) is the same
functional with the evaluator samples added to \(Y\) and \(\kappa\) the
physical hidden conditional.
\status{ESTABLISHED}

\noindent\emph{Proof.}
All identities below are the chain rule for relative entropy on standard
Borel spaces and hold in \([0,\infty]\). For (i), let \(J\) be the joint law
of \((B,Y,Z)\), and let \(\widehat J\) keep the joint law of \((B,Y)\) and
draw \(Z\) from \(\kappa_Y\). Disintegrating in \((B,Y)\) gives
\(\KL(J\Vert\widehat J)=\sum_bw_b\mathcal H_\kappa(Q_b)\), because the
conditional law of \(Z\) given \((B=b,Y)\) is the hidden conditional of
\(Q_b\). The \((Y,Z)\)-marginal of \(\widehat J\) is
\(\operatorname{Law}(Y)\otimes\kappa\), and under \(\widehat J\) the variable
\(B\) is conditionally independent of \(Z\) given \(Y\). Disintegrating in
\((Y,Z)\) therefore gives
\(\KL(J\Vert\widehat J)=\mathcal H_\kappa(Q)
+\E\KL(\operatorname{Law}(B\mid Y,Z)\Vert\operatorname{Law}(B\mid Y))\), and
the last term is \(\mathbb I(B;Z\mid Y)\), the difference of the conditional
Shannon entropies of \(B\) given \(Y\) and given \((Y,Z)\); it lies between
zero and the entropy of \(B\). Convexity follows because the information term
is nonnegative. For an event \(A\) with \(r=\Pr(A^c)\), the case \(r=0\) is
trivial; otherwise \(B=\mathbf 1_A\) gives
\((1-r)\mathcal H_\kappa(Q_A)\le\mathcal H_\kappa(Q)+\mathcal S_2(r)\), and
\(\mathcal S_2\) increases on \([0,1/2]\) while \(1-r\ge1-\bar r\). The
same identity with the term of \(A^c\) dropped gives
\(\mathcal H_\kappa(Q)\ge\Pr(A)\mathcal H_\kappa(Q_A)-\mathcal S_2(\Pr(A))\).

For (ii), the chain rule gives
\(\KL(\operatorname{Law}(Y,\Upsilon,Z)\Vert\operatorname{Law}(Y,\Upsilon)\otimes\kappa)
=\E\KL(\operatorname{Law}(Z\mid Y,\Upsilon)\Vert\kappa_Y)\). Forgetting
\(\Upsilon\) maps the two laws to \(\operatorname{Law}(Y,Z)\) and
\(\operatorname{Law}(Y)\otimes\kappa\), whose relative entropy is
\(\mathcal H_\kappa(Q)\), and marginalization does not increase relative
entropy.

For (iii), fix \(y\) and put \(P=\operatorname{Law}(Z\mid Y=y)\),
\(R=\kappa_y\), and \(\mathrm g_n=\mathrm g(y,\cdot)\wedge n\). It suffices
to show \(\KL(P\Vert R)\ge\int\mathrm g(y,z)P(dz)-\log\mathrm c_1\) and to
average over \(Y\). The case \(\KL(P\Vert R)=\infty\) is trivial, so let
\(\KL(P\Vert R)<\infty\). With
\(R_n=e^{\mathrm g_n}R/\E_Re^{\mathrm g_n}\),
\(\log(dP/dR)=\log(dP/dR_n)+\mathrm g_n-\log\E_Re^{\mathrm g_n}\), and
integration against \(P\) gives
\(\KL(P\Vert R)\ge\E_P\mathrm g_n-\log\E_Re^{\mathrm g_n}
\ge\E_P\mathrm g_n-\log\mathrm c_1\), and monotone convergence as
\(n\to\infty\) proves (iii).

On a matching stratum \(\mathbf x=y+F_Sz\) is a linear isometry between
\(\R^{12}\) and \(W_S\times\R^3\), so conditioning on \(C(\mathbf X,S)\) is
conditioning on \(Y\), and off the matchings the reference is the point mass
at the constant \(Z\); this identifies \(\mathcal H_{\mathrm{opt}}\). The
physical functional is identified in the same way with the retained variables
\((C(\mathbf X,S),U)\). \(\square\)

\lemmaheading{Physical and optimizing discarded information}{lem:phil-early-physical-comparison}
In the four-cycle at any strictly positive rows, let \(q\) be a law of
\((\mathbf X,S)\) with
\(m_{\mathrm h}(q)=\E_q[\mathbf 1_{\{S^{(0)},S^{(1)}\}}(S)|F_S^\top\mathbf X|^2]<\infty\).
Then
\[
 \mathcal H_{\mathrm{opt}}(q)-4m_{\mathrm h}(q)-2\mathrm T
 \le\mathcal H_{\mathrm{phys}}(q)
 \le\mathcal H_{\mathrm{opt}}(q)+\log\frac{4+5\mathrm T}{4+\mathrm T}
 \le\mathcal H_{\mathrm{opt}}(q)+\mathrm T.
\]
\status{ESTABLISHED}

\noindent\emph{Proof.}
Off the matchings neither functional discards anything. Fix a matching \(M\)
and a retained point \(y\), write \(q_y=q(dz\mid M,y)\), \(v=1/2\),
\(d_y=\mathfrak Dy\), \(d_z=\mathfrak DF_Mz\), and let
\(\Pi_{\mathrm T,Z}\) and the probability law \(\overline P_y\) be those of
\cref{eq:phil-conditioned-evaluator-tilt}. The proof of
\cref{prop:phil-actual-conditioned-parent} shows that \(\overline P_y\) has
\(z\)-marginal \(\Pi_{\mathrm T,Z}\), that given \(z\) its evaluator law is
\(\mathcal N(\bar m(y+F_Mz),\sigma^2_{\mathrm{ev}}I)\) for \(m\) times
\(\mathcal N(d_y,vI)\) for \(k-m\), that its conditional law of \(z\) given the evaluator sample \(U\) is
the physical hidden conditional, and that
\(\KL(q_y\Vert\Pi_{\mathrm T,Z})+\E_{q_y}|d_z|^2/(2v)
=\KL(q_y\Vert\nu_{\mathrm T,Z})+\log[(4+5\mathrm T)/(4+\mathrm T)]\). The
recognition lift \(Q_y(dz,dU)=q_y(dz)\nu^*_{y+F_Mz}(dU)\) of
\cref{eq:rg-formation-evaluator-lift} has, given \(z\), the same law of \(m\)
and increment law \(\mathcal N(d_y+d_z,vI)\), at relative entropy
\(|d_z|^2/(2v)\). Computing \(\KL(Q_y\Vert\overline P_y)\) through \(z\) and
through \(U\) gives, in \([0,\infty]\),
\[
 \mathcal H_{\mathrm{phys}}(q\mid M,y)
 +\KL\bigl(Q_y(dU)\Vert\overline P_y(dU)\bigr)
 =\KL(q_y\Vert\nu_{\mathrm T,Z})+\log\frac{4+5\mathrm T}{4+\mathrm T}.
\]
Dropping the second term on the left gives the upper bound. For the lower
bound, the two
evaluator marginals are the \(U\)-marginals of
\(q_y(dz)\Pi_{\mathrm T,Z}(dz')\nu^*_{y+F_Mz}(dU)\) and
\(q_y(dz)\Pi_{\mathrm T,Z}(dz')\overline P_y(dU\mid z')\), so marginalization
gives
\[
 \KL\bigl(Q_y(dU)\Vert\overline P_y(dU)\bigr)
 \le\iint\bigl(|\bar m(F_M(z-z'))|^2+|d_z|^2\bigr)
     q_y(dz)\Pi_{\mathrm T,Z}(dz'),
\]
the integrand being the relative entropy of two Gaussian laws with common
covariance \(I/2\), since \(\sigma^2_{\mathrm{ev}}=v=1/2\). Since \(\bar m\) is a block of coordinates, \(F_M\) has
orthonormal columns, and \(\mathfrak D\mathfrak D^\top=2I_4\), the integrand
is at most \(|z-z'|^2+2|z|^2\le4|z|^2+2|z'|^2\), and
\(\E_{\Pi_{\mathrm T,Z}}|z'|^2=\mathrm T[1/(8+\mathrm T)+2/(4+\mathrm T)]\le5\mathrm T/8\).
Since \(m_{\mathrm h}(q)<\infty\), \(\E_{q_y}|z|^2<\infty\) for almost every
retained point \(y\), and for such \(y\) the relative entropy of the
evaluator marginals is at most \(4\E_{q_y}|z|^2+5\mathrm T/4<\infty\);
subtracting it from the identity
gives
\(\mathcal H_{\mathrm{phys}}(q\mid M,y)\ge\KL(q_y\Vert\nu_{\mathrm T,Z})
-4\E_{q_y}|z|^2-5\mathrm T/4\). Integrate over the retained law, whose
matching strata have total mass at most one. Finally
\(\log[(4+5\mathrm T)/(4+\mathrm T)]\le4\mathrm T/(4+\mathrm T)\le\mathrm T\).
\(\square\)

\lemmaheading{Rows and the held seed}{lem:phil-early-held-seed}
In \cref{def:phil-early-seeded-prefix} the following hold.

(i) For every \(\epsilon_{\mathrm L}\ge0\) the rows exist uniquely, are
deterministic and continuous, and every entry lies in
\([e^{-\epsilon_{\mathrm L}t}/2,1-e^{-\epsilon_{\mathrm L}t}/2]\). The two
entries of a row have ratio at most \(\Lambda_{\mathrm L}(t)\); a receiver
\(i\) with sources \(j,k\) satisfies
\(\beta_{ik}(\mathbf x;\pi(t))\le\Lambda_{\mathrm L}(t)
\exp[-(\mathcal D_{e(i,k)}(\mathbf x)-\mathcal D_{e(i,j)}(\mathbf x))/\mathrm T]\);
and each of its two choice probabilities is at most
\(\Lambda_{\mathrm L}(t)\) times its value at rows \(1/2\).

(ii) For each matching \(M\), \(F_M^\top\mathbf x^\circ=0\) and
\(z_+(\mathbf x^\circ)=-1\). Along \(\bar{\mathbf x}^{\sigma,M}\) the hidden
coordinates vanish, \(z_+=-\sigma e^{-(8+\mathrm T)t}\), every selected
residual equals \(-\sigma e^{-(8+\mathrm T)t}\), and every exterior residual
equals \(\sigma e^{-\mathrm Tt}(-3/2+e^{-8t}/2)\), in both channels.

(iii) Suppose \(|\mathbf X_0-\sigma\mathbf x^\circ|<\varrho_\circ\) and
\(\sup_{s\le t}|\mathbf Z^M_s|\le\delta_\epsilon\), and let \(u\le t\). Then
\(\widetilde{\mathbf x}=\widetilde{\mathbf X}^M_u\) satisfies
\(|\widetilde{\mathbf x}-\bar{\mathbf x}^{\sigma,M}(u)|\le\varrho_\circ+\delta_\epsilon\)
and \(|F_M^\top\widetilde{\mathbf x}|\le\varrho_\circ+\delta_\epsilon\); every
selected residual satisfies
\(|r(\widetilde{\mathbf x})|\le e^{-(8+\mathrm T)u}+\sqrt3\varrho_\circ e^{-4u}+\epsilon^2/8\),
and every exterior residual satisfies
\(|r(\widetilde{\mathbf x})|\ge e^{-\mathrm Tu}(3/2-e^{-8u}/2)-\sqrt3(\varrho_\circ+\delta_\epsilon)\).
At \(u=0\) the internal belief relative entropy of each matched pair exceeds
\(1/4\). If \(\mathrm Tt\le1/12\) and \(u\in[1/4,t]\), each receiver's
exterior cost exceeds its partner cost by more than one and every exterior
belief relative entropy exceeds \(1/2\). If \(u\ge\frac12\log(1/\epsilon)\),
every internal belief and model relative entropy of the two matched pairs is
at most \(\epsilon^4\).

(iv) Let \(s\ge0\), let \(\mathbf X_s\) be independent of
\((W_{s+r}-W_s)_{r\ge0}\) with \(|F_M^\top\mathbf X_s|\le\mathrm R\), and let
\(\widehat{\mathbf X}_{s+r}=e^{-K_Mr}\mathbf X_s
+\sqrt{2\mathrm T}\int_s^{s+r}e^{-K_M(s+r-v)}dW_v\). For \(r>0\),
\[
 \mathcal H_{\mathrm{opt}}\bigl(\operatorname{Law}(\widehat{\mathbf X}_{s+r},M)\bigr)
 \le\frac{9\mathrm R^2e^{-8r}}{2\mathrm T}+\frac{3e^{-16r}}{4(1-e^{-8r})}.
\]
\status{ESTABLISHED}

\noindent\emph{Proof.}
(i) The proof of \cref{prop:phil-seeded-contrast-formation}(a) uses the
mobility only as a positive constant; with \(\epsilon_{\mathrm L}\) in place
of \(\epsilon\) it gives existence, uniqueness, continuity, and
\(|\pi_{ij}(t)-1/2|\le(1-e^{-\epsilon_{\mathrm L}t})/2\), and the case
\(\epsilon_{\mathrm L}=0\) is trivial. The ratio of the endpoints of the
entry interval is \(\Lambda_{\mathrm L}(t)\). The bound on \(\beta_{ik}\)
drops the partner term from its denominator. With likelihood factors
\(\mathrm c_1,\mathrm c_2>0\) and entry \(\mathrm p\), the ratio of
\(\mathrm p\mathrm c_1/(\mathrm p\mathrm c_1+(1-\mathrm p)\mathrm c_2)\) to
its value at \(\mathrm p=1/2\) is
\(\mathrm p(\mathrm c_1+\mathrm c_2)/(\mathrm p\mathrm c_1+(1-\mathrm p)\mathrm c_2)
\le\max\{1,\mathrm p/(1-\mathrm p)\}\), and similarly for the other choice.

(ii) Every residual of \(\mathbf x^\circ\) equals \(-1\). Since
\(\mathbf a^b_g=\sqrt2n^+_g+n^-_g\) and \(\mathbf a^m_g=\sqrt2n^+_g-n^-_g\),
this gives \(n^-_g\cdot\mathbf x^\circ=0\) and
\(\sqrt2n^+_g\cdot\mathbf x^\circ=-1\) for \(g\in\{e,f\}\), hence
\(F_M^\top\mathbf x^\circ=0\) and
\(z_+(\mathbf x^\circ)=(\sqrt2n^+_e+\sqrt2n^+_f)\cdot\mathbf x^\circ/2=-1\).
The residuals along the center path are
\cref{eq:rg-formation-cycle-center} for \(S^{(0)}\); the rotation
\(i\mapsto i+1\) fixes \(\mathbf x^\circ\) and maps \(S^{(0)}\) to
\(S^{(1)}\), and the sign \(\sigma\) enters by linearity. Since
\(K_MF_M=F_MK_Z\) and \(u_M\) is an eigenvector of \(K_M\) with eigenvalue
\(8+\mathrm T\), the hidden coordinates stay zero and
\(z_+=-\sigma e^{-(8+\mathrm T)t}\).

(iii) Since \(K_M\succeq\mathrm TI\), \(e^{-K_Mu}\) is a contraction, so
\(|\widetilde{\mathbf x}-\bar{\mathbf x}^{\sigma,M}(u)|
\le|\mathbf X_0-\sigma\mathbf x^\circ|+|\mathbf Z^M_u|\), and the hidden
coordinates of the center path vanish. For a residual covector \(\mathbf a\)
of norm \(\sqrt3\),
\(r(\widetilde{\mathbf x})=r(\bar{\mathbf x}^{\sigma,M}(u))
+(e^{-K_Mu}\mathbf a)\cdot(\mathbf X_0-\sigma\mathbf x^\circ)
+\mathbf a\cdot\mathbf Z^M_u\). A selected covector lies in the span of
eigenvectors of \(K_M\) with eigenvalues at least \(4\), so
\(|e^{-K_Mu}\mathbf a|\le\sqrt3e^{-4u}\); an exterior one has
\(|e^{-K_Mu}\mathbf a|\le\sqrt3\); and \(\sqrt3\delta_\epsilon=\epsilon^2/8\).
With (ii) this gives the two residual bounds. At \(u=0\) a selected residual
has magnitude at least \(1-\sqrt3\varrho_\circ\), and
\((1-\sqrt3\varrho_\circ)^2/2>0.48\). For \(u\in[1/4,t]\) with
\(\mathrm Tt\le1/12\), the selected residuals are at most
\(e^{-2}+\sqrt3(\varrho_\circ+\delta_\epsilon)<0.153<1/5\), and the exterior
ones at least
\(e^{-1/12}(3/2-e^{-2}/2)-\sqrt3(\varrho_\circ+\delta_\epsilon)>1.3004>13/10\);
the partner cost \([(r^b)^2+(r^m)^2]/2\) is below \(1/25\), the exterior cost
above \(169/100\), and every exterior belief relative entropy
\((r^b)^2/2\) exceeds \(169/200\). For \(u\ge\frac12\log(1/\epsilon)\) a
selected residual is at most
\(\epsilon^4+\sqrt3\varrho_\circ\epsilon^2+\epsilon^2/8<0.15\epsilon^2\),
so the belief relative entropy \((r^b)^2/2\) and the model relative entropy
\((r^m)^2/(2\sigma^2_{\mathrm{ev}})=(r^m)^2\) are at most \(\epsilon^4\).

(iv) The matrix \(K_M\) preserves \(W_M\) and \(Z_M\), and the projections of
the increments of \(W\) onto these orthogonal subspaces are independent
Brownian motions. Hence the retained part of
\(\widehat{\mathbf X}_{s+r}\) is a function of \(\mathbf X_s\) and the
retained increments, while its hidden part is
\(e^{-K_Zr}z_s+\xi_r\), with \(z_s=F_M^\top\mathbf X_s\) and
\(\xi_r\sim\mathcal N(0,\mathrm TK_Z^{-1}(I-e^{-2K_Zr}))\) independent of
\(\mathbf X_s\) and of the retained increments. By
\cref{lem:phil-early-discarded-information}(ii) with \(\Upsilon=\mathbf X_s\), the
left side is at most the mean relative entropy of
\(\mathcal N(e^{-K_Zr}z_s,\mathrm TK_Z^{-1}(I-e^{-2K_Zr}))\) to
\(\nu_{\mathrm T,Z}\). For a coordinate of rate \(\lambda_k\) and
\(a_k=e^{-2\lambda_kr}\), this relative entropy is
\(\frac12[-a_k-\log(1-a_k)+\lambda_ka_kz_{s,k}^2/\mathrm T]\). The rates lie
in \([4,9]\), so \(\lambda_ka_k\le9e^{-8r}\) and \(a_k\le e^{-8r}\); moreover
\(|z_s|\le\mathrm R\) and \(-a-\log(1-a)=\sum_{n\ge2}a^n/n\le a^2/(2(1-a))\).
\(\square\)

\propositionheading{The seeded early prefix at a deterministic activation time}{prop:phil-early-seeded-parent}
In \cref{def:phil-early-seeded-prefix}, let \(\theta\ge1/4\) with
\(\mathrm T\theta\le1/12\). Put
\[
 \mathsf n_\theta=\min\Bigl\{73728\epsilon^4\theta,
   24\exp\Bigl[-\frac{\epsilon^{-4}}{36864\theta}\Bigr]\Bigr\},\qquad
 p^\circ_\theta=\mathsf n_\theta
   +4\Lambda_{\mathrm L}(\theta)\bigl(\theta-\tfrac14\bigr)e^{-1/\mathrm T},
\]
\[
 \eta^\circ(\theta)=\frac{9\varrho_\circ^2e^{-8\theta}}{2\mathrm T}
   +\frac{3e^{-16\theta}}{4(1-e^{-8\theta})},\qquad
 \bar\eta^\circ(\theta)
   =\frac{\eta^\circ(\theta)+\mathcal S_2(p^\circ_\theta)}{1-p^\circ_\theta},
 \qquad
 \mathrm c_{\mathsf m}=e^{-2\varrho_\circ(|\mathsf m_0|+2+\varrho_\circ)},
\]
and assume \(p^\circ_\theta<1/2\). Then the following hold.

(a) \emph{Prefix.} The four seeds are disjoint,
\(\Pr((\mathscr A^M_\theta)^c\mid\mathscr B_{\sigma,M})\le p^\circ_\theta\)
for each of them, and
\(\Pr(\mathscr E^\circ_\theta)\ge(1-p^\circ_\theta)\Pr(\mathscr B^\circ)
\ge\mathsf c(\mathsf m_0)\), with \(\mathsf c\) from
\cref{prop:phil-early-block-null}. On
\(\mathscr B_{\sigma,M}\cap\mathscr A^M_\theta\) the native path equals
\(\widetilde{\mathbf X}^M\) on \([0,\theta]\), \(S_\theta=M\), and
\(|\mathbf X_\theta-\bar{\mathbf x}^{\sigma,M}(\theta)|\le\varrho_\circ+\delta_\epsilon\);
in particular \(q^\circ_\theta\) is carried by the matching strata with hidden
coordinates of norm at most \(\varrho_\circ+\delta_\epsilon\).

(b) \emph{Upper bound.}
\[
 \mathcal H_{\mathrm{opt}}(q^\circ_\theta)\le\bar\eta^\circ(\theta),
 \qquad
 \mathcal H_{\mathrm{phys}}(q^\circ_\theta)
 \le\bar\eta^\circ(\theta)+\log\frac{4+5\mathrm T}{4+\mathrm T}
 \le\bar\eta^\circ(\theta)+\mathrm T .
\]
The rows enter only through \(\Lambda_{\mathrm L}(\theta)\), and the bound
does not depend on \(\mathsf m_0\). For \(\epsilon_{\mathrm L}\in[0,1]\), as
\(\epsilon\to0\) with \(\epsilon^4\theta\to0\) and
\(\theta-\log(1/\epsilon)\) bounded below,
\(\bar\eta^\circ(\theta)=\frac92\varrho_\circ^2e^{-8(\theta-\log(1/\epsilon))}+o(1)\),
which tends to zero when \(\theta-\log(1/\epsilon)\to\infty\); uniformly on
\([t_\epsilon,t^\circ_*]\),
\(\bar\eta^\circ(\theta)=\frac92\varrho_\circ^2e^{-8(\theta-t_\epsilon)}\mathrm T
+O(\mathrm T^4)\).

(c) \emph{Lower bound.} If \(\theta\le\log(1/\epsilon)\), then
\[
 \mathcal H_{\mathrm{opt}}(q^\circ_\theta)
 \ge\frac{4+5\mathrm T}{4\mathrm T}
 \Bigl[\frac{\varrho_\circ^2}7\mathrm c_{\mathsf m}e^{-2(4+5\mathrm T)\theta}
       -3\varrho_\circ^2(p^\circ_\theta)^{1/2}\Bigr]-\log2,
\]
and \(\mathcal H_{\mathrm{phys}}(q^\circ_\theta)\ge
\mathcal H_{\mathrm{opt}}(q^\circ_\theta)-1/2000\). For each fixed
\(\mathsf m_0\) and \(\epsilon_{\mathrm L}\in[0,1]\), both tend to infinity
as \(\epsilon\to0\) with \(\log(1/\epsilon)-\theta\to\infty\), and at
\(\theta=\theta_\circ\) they are at least
\((e^{-8}\varrho_\circ^2/7)\mathrm c_{\mathsf m}\mathrm T^{-1/2}(1+o(1))\).

(d) \emph{Future marginals under the same prefix.} For \(t\ge\theta\) with
\(\mathrm Tt\le1/12\) and \(p^\circ_t<1/2\), put
\(\mathfrak L_t=17(|\mathsf m_0|^2+14)/(2\mathrm T)
+4(\log2+\epsilon_{\mathrm L}t)(t+t^2)^{1/2}\). Then
\[
 \mathcal H_{\mathrm{opt}}\bigl(\operatorname{Law}(\mathbf X_t,S_t\mid
 \mathscr E^\circ_\theta)\bigr)
 \le\bar\eta^\circ(t)+\frac{p^\circ_t\log\frac1{p^\circ_t\Pr(\mathscr B^\circ)}
 +\bigl(p^\circ_t/\Pr(\mathscr B^\circ)\bigr)^{1/2}\mathfrak L_t}{1-p^\circ_\theta},
\]
and \(\mathcal H_{\mathrm{phys}}\) of the same law is at most the right side
plus \(\mathrm T\). For each fixed \(\mathsf m_0\) and
\(\epsilon_{\mathrm L}\in[0,1]\), the fraction tends to zero as
\(\epsilon\to0\), uniformly for
\(\theta\le t\le\epsilon^{-4}/(663552\log(1/\epsilon))\).

(e) \emph{A path-measurable prefix.} For \(\theta>1/4\), let
\(\mathscr E^{\mathrm{path}}_\theta\) be the union over seeds of the events
\(\{|\mathbf X_0-\sigma\mathbf x^\circ|<\varrho_\circ\}\cap\{S_s=M\text{ for }s\le\theta\}
\cap\{\sup_{s\le\theta}|\mathbf X_s-e^{-K_Ms}\mathbf X_0|\le\delta_\epsilon\}\),
a function of the configuration and assignment path on \([0,\theta]\) that
contains \(\mathscr E^\circ_\theta\). Put
\(p^{\mathrm{path}}_\theta=e^{1/4}\mathsf n_\theta/(1-\mathsf n_\theta)
+4\Lambda_{\mathrm L}(\theta)(\theta-\frac14)e^{-1/\mathrm T}\), and let
\(\eta^{\mathrm{path}}(\theta)\) be \(\eta^\circ(\theta-\frac14)\) with
\(\varrho_\circ\) replaced by \(\varrho_\circ+\delta_\epsilon\). If
\(p^{\mathrm{path}}_\theta<1/2\), then
\(\mathcal H_{\mathrm{opt}}(\operatorname{Law}(\mathbf X_\theta,S_\theta\mid
\mathscr E^{\mathrm{path}}_\theta))\le
[\eta^{\mathrm{path}}(\theta)+\mathcal S_2(p^{\mathrm{path}}_\theta)]
/(1-p^{\mathrm{path}}_\theta)\), and \(\mathcal H_{\mathrm{phys}}\) of that
law is at most this bound plus \(\mathrm T\).
\status{ESTABLISHED}

Parts (b) and (c) fix the activation clock up to the window
\(\theta=\log(1/\epsilon)+O(1)\). On that window both bounds are of order one,
and the estimates leave open whether the discarded information tends to zero
there. The window \([t_\epsilon,t^\circ_*]\) of
\cref{prop:rg-formation-endogenous-cycle} lies \(\log(1/\epsilon)\) beyond the
threshold. Only the discarded conditional term is bounded; the relative entropy
of \(q^\circ_\theta\) to the restricted optimizer also contains a retained
term, which is not bounded here.

\noindent\emph{Proof.}
(a) The balls about \(\pm\mathbf x^\circ\) are disjoint because
\(|\mathbf x^\circ|=2>\varrho_\circ\), and \(S_0\) separates the matchings.
The initial state and the clock are independent, so
\(\Pr(\mathscr B_{\sigma,M})=\frac{e^{-1/4}}{16}\mathcal N(\mathsf m_0,I_{12})
\{|\mathbf x-\sigma\mathbf x^\circ|<\varrho_\circ\}\), and the proof of
\cref{prop:phil-seeded-contrast-formation}(b) bounds this below by
\(\mathsf c(\mathsf m_0)\) for \(\sigma=1\); hence
\(\Pr(\mathscr B^\circ)\ge2\mathsf c(\mathsf m_0)\). On
\(\mathscr B_{\sigma,M}\cap\{S_s=M\text{ for }s\le t\}\) the native configuration and
\(\widetilde{\mathbf X}^M\) solve the same linear equation with the same
Brownian path between rings, and a ring that keeps the assignment does not
move the configuration, so they agree on \([0,t]\); with
\cref{lem:phil-early-held-seed}(iii) this gives the last assertions.

The tube event depends only on \(W\), which is independent of the seed.
Doob's \(L^2\) maximal inequality gives
\(\E\sup_{s\le\theta}|\mathbf Z^M_s|^2\le384\mathrm T\theta\), as in the
paragraph \emph{Episode} of the proof of
\cref{prop:rg-formation-endogenous-cycle}, and Chebyshev's inequality bounds
the tube failure by \(384\mathrm T\theta/\delta_\epsilon^2=73728\epsilon^4\theta\).
In an orthonormal eigenbasis of \(K_M\) the same paragraph gives
\(\sup_{s\le\theta}|Z^k_s|\le2\sqrt{2\mathrm T}\sup_{s\le\theta}|W^k_s|\), so
the tube fails only if some \(|W^k|\) exceeds
\(\mathrm a_\delta=\delta_\epsilon/(4\sqrt{6\mathrm T})\) before \(\theta\), where
\(\mathrm a_\delta^2=\epsilon^{-4}/18432\). The reflection principle and
\(\Pr(\mathcal N(0,1)\ge x)\le e^{-x^2/2}/2\) bound this by
\(2e^{-\mathrm a_\delta^2/(2\theta)}\) for each of the twelve coordinates. Hence the
tube fails with probability at most \(\mathsf n_\theta\).

For departures, let \(N\) count the rings in \((1/4,\theta]\) that are
preceded by the good state, in which the seed holds and the assignment has
equaled \(M\) and the noise has stayed in the tube up to the ring, and at
which the redraw differs from \(M\). In the good state the configuration is
held, so by \cref{lem:phil-early-held-seed}(iii) every exterior cost exceeds
the partner cost by more than one, and by part (i) of that lemma the redraw
leaves \(M\) with conditional probability at most
\(4\Lambda_{\mathrm L}(\theta)e^{-1/\mathrm T}\), the variables of the ring
being independent of the past. The seed is \(\mathcal F_{1/4}\)-measurable,
so the compensator of the clock gives
\(\E[N\mid\mathscr B_{\sigma,M}]\le4\Lambda_{\mathrm L}(\theta)
(\theta-\frac14)e^{-1/\mathrm T}\). If the tube holds on \([0,\theta]\) and
the assignment leaves \(M\) by time \(\theta\), the first departure occurs
after \(1/4\) from the good state, so \(N\ge1\); Markov's inequality completes
the remainder bound. Finally
\(\Pr(\mathscr E^\circ_\theta)=\sum\Pr(\mathscr B_{\sigma,M})
\Pr(\mathscr A^M_\theta\mid\mathscr B_{\sigma,M})\ge(1-p^\circ_\theta)
\Pr(\mathscr B^\circ)\), which is at least \(\mathsf c(\mathsf m_0)\) because
\(p^\circ_\theta<1/2\).

(b) Fix a seed and write \(\mathscr B=\mathscr B_{\sigma,M}\),
\(\mathscr A=\mathscr A^M_\theta\). Under \(\Pr(\cdot\mid\mathscr B)\) the
Brownian motion keeps its law and is independent of \(\mathbf X_0\), and
\(|F_M^\top\mathbf X_0|=|F_M^\top(\mathbf X_0-\sigma\mathbf x^\circ)|<\varrho_\circ\)
by \cref{lem:phil-early-held-seed}(ii). By part (iv) of that lemma with
\(s=0\) and \(\mathrm R=\varrho_\circ\), the law \(\widehat Q\) of
\((\widetilde{\mathbf X}^M_\theta,M)\) under \(\Pr(\cdot\mid\mathscr B)\)
has \(\mathcal H_{\mathrm{opt}}(\widehat Q)\le\eta^\circ(\theta)\). By (a),
the law of \((\mathbf X_\theta,S_\theta)\) given
\(\mathscr B\cap\mathscr A\) is the conditional law of \(\widehat Q\) given
\(\mathscr A\), and \(\Pr(\mathscr A^c\mid\mathscr B)\le p^\circ_\theta<1/2\),
so \cref{lem:phil-early-discarded-information}(i) bounds its discarded
information by \(\bar\eta^\circ(\theta)\). The law \(q^\circ_\theta\) is the
mixture of these four seed laws, and convexity gives the first bound. The
physical bound follows from \cref{lem:phil-early-physical-comparison}, the
hidden coordinates being bounded by (a).

For the asymptotic statements, \(\Lambda_{\mathrm L}(\theta)\le2e^\theta\)
when \(\epsilon_{\mathrm L}\le1\), so
\(p^\circ_\theta\le24\exp[-\epsilon^{-4}/(36864\theta)]
+8\theta e^{\theta}e^{-1/\mathrm T}\). If \(\epsilon^4\theta\to0\), then
\(\theta\le\epsilon^{-4}\) eventually and both terms tend to zero, hence so
do \(\mathcal S_2(p^\circ_\theta)\) and
\(p^\circ_\theta\eta^\circ(\theta)\) whenever \(\eta^\circ(\theta)\) stays
bounded. Moreover
\(9\varrho_\circ^2e^{-8\theta}/(2\mathrm T)
=\frac92\varrho_\circ^2e^{-8(\theta-\log(1/\epsilon))}\), and the second
term of \(\eta^\circ(\theta)\) tends to zero because \(\theta\to\infty\). On
\([t_\epsilon,t^\circ_*]\) one has \(e^{-8\theta}=\mathrm T^2
e^{-8(\theta-t_\epsilon)}\), so the second term is at most
\(\mathrm T^4/(1-\mathrm T^2)\), and
\(\epsilon^{-4}/(36864t^\circ_*)\) grows faster than every multiple of
\(\log(1/\epsilon)\), so \(p^\circ_\theta\) and
\(\mathcal S_2(p^\circ_\theta)\) are \(O(\mathrm T^k)\) for every \(k\).

(c) Put \(z_\perp=(z_2,z_3)\) on each matching stratum. Under
\(\nu_{\mathrm T,Z}\) the coordinates \(z_2,z_3\) are independent with
variance \(\mathrm T/(4+5\mathrm T)\), so
\(\mathrm g=(4+5\mathrm T)|z_\perp|^2/(4\mathrm T)\) has
\(\int e^{\mathrm g}d\nu_{\mathrm T,Z}=2\), and
\cref{lem:phil-early-discarded-information}(iii) gives
\(\mathcal H_{\mathrm{opt}}(q)\ge\frac{4+5\mathrm T}{4\mathrm T}
\E_q|z_\perp|^2-\log2\) for every law \(q\) carried by the matching strata.
Since \(q^\circ_\theta\) is a mixture of the seed laws and
\(\Pr(\mathscr A\mid\mathscr B)\le1\),
\(\E_{q^\circ_\theta}|z_\perp|^2\) is at least the least over seeds of
\(\E[|\widetilde z_\perp(\theta)|^2\mathbf 1_{\mathscr A}\mid\mathscr B]\),
where \(\widetilde z\) are the hidden coordinates of the held process. In
these, \(\widetilde z_\perp(\theta)=e^{-(4+5\mathrm T)\theta}z_\perp(\mathbf X_0)
+\xi_\perp\), with \(\xi\) centered and independent of \(\mathbf X_0\) under
\(\Pr(\cdot\mid\mathscr B)\), so
\(\E[|\widetilde z_\perp(\theta)|^2\mid\mathscr B]\ge
e^{-2(4+5\mathrm T)\theta}\E[|z_\perp(\mathbf X_0)|^2\mid\mathscr B]\). Given
the seed, the offset \(\mathbf X_0-\sigma\mathbf x^\circ\) has a density with
respect to the uniform law on the ball of radius \(\varrho_\circ\) that is at
least \(\mathrm c_{\mathsf m}\): the squared distances from \(\mathsf m_0\)
of two points of the ball differ by at most
\(4\varrho_\circ(|\sigma\mathbf x^\circ-\mathsf m_0|+\varrho_\circ)\), so the
Gaussian density on the ball varies by at most the factor
\(\exp[2\varrho_\circ(|\sigma\mathbf x^\circ-\mathsf m_0|+\varrho_\circ)]\),
and \(|\sigma\mathbf x^\circ-\mathsf m_0|\le|\mathsf m_0|+2\). Under the uniform law each orthonormal coordinate has
second moment \(\varrho_\circ^2/14\), and
\(z_\perp(\mathbf X_0)=z_\perp(\mathbf X_0-\sigma\mathbf x^\circ)\) consists
of two of them, so
\(\E[|z_\perp(\mathbf X_0)|^2\mid\mathscr B]\ge
\mathrm c_{\mathsf m}\varrho_\circ^2/7\). On the other side
\(|\widetilde z(\theta)|\le\varrho_\circ+|\xi|\) and
\(\operatorname{Cov}\xi\preceq(\mathrm T/4)I_3\), so
\(\E[|\widetilde z(\theta)|^4\mid\mathscr B]\le8\varrho_\circ^4
+8\cdot\frac{15}{16}\mathrm T^2\le9\varrho_\circ^4\), and the Cauchy--Schwarz
inequality with (a) gives
\(\E[|\widetilde z_\perp(\theta)|^2\mathbf 1_{\mathscr A^c}\mid\mathscr B]\le
3\varrho_\circ^2(p^\circ_\theta)^{1/2}\). This proves the display. The
physical bound is \cref{lem:phil-early-physical-comparison} with
\(4(\varrho_\circ+\delta_\epsilon)^2+2\mathrm T<1/2000\).

For the limits, write \(\theta=\log(1/\epsilon)-\mathrm s\). Then
\(e^{-2(4+5\mathrm T)\theta}/\mathrm T=e^{8\mathrm s}e^{-10\mathrm T\theta}\)
with \(\mathrm T\theta\to0\), while
\(p^\circ_\theta\le24\exp[-\epsilon^{-4}/(36864\log(1/\epsilon))]
+8\epsilon^{-1}\log(1/\epsilon)e^{-1/\mathrm T}\), so
\((p^\circ_\theta)^{1/2}/\mathrm T\to0\). Hence the lower bound is
\(\frac{\varrho_\circ^2}7\mathrm c_{\mathsf m}e^{8\mathrm s}(1+o(1))-\log2\),
which tends to infinity with \(\mathrm s\); at \(\theta=\theta_\circ\),
\(e^{8\mathrm s}=e^{-8}\mathrm T^{-1/2}\).

(d) Decompose \(\mathscr E^\circ_\theta\) into \(\mathscr E^\circ_t\) and
\(\mathscr G_t=\mathscr E^\circ_\theta\setminus\mathscr E^\circ_t\), since
\(\mathscr A^M_t\subseteq\mathscr A^M_\theta\). Convexity bounds the left side
by \([\Pr(\mathscr E^\circ_t)\bar\eta^\circ(t)+\Pr(\mathscr G_t)
\mathcal H_{\mathrm{opt}}(q^{\mathscr G_t}_t)]/\Pr(\mathscr E^\circ_\theta)\),
by (b) at time \(t\), where \(q^{\mathscr G_t}_t\) is the law of
\((\mathbf X_t,S_t)\) given \(\mathscr G_t\). Since the two seed balls are
disjoint, \(\Pr(\mathscr B^\circ)\le e^{-1/4}/8\), so (a) gives
\(\Pr(\mathscr G_t)\le p^\circ_t\Pr(\mathscr B^\circ)\le1/16\) and
\(\Pr(\mathscr E^\circ_\theta)\ge(1-p^\circ_\theta)\Pr(\mathscr B^\circ)\).

To bound \(\mathcal H_{\mathrm{opt}}(q^{\mathscr G_t}_t)\), let
\(\mathcal F^{\mathrm{ch}}_t\) be generated by the initial state, by \(W\) and the clock on
\([0,t]\), and by the source choices at the rings in \([0,t]\); the
configuration and assignment path on \([0,t]\) is a measurable function of
these data, and \(\mathscr G_t\in\mathcal F^{\mathrm{ch}}_t\). Let \(\mathsf Q\) be the law
of the same data for the kinetics with rows frozen at \(1/2\) started from
\(\mu_{\mathrm T,1/2}\), which those kinetics leave invariant
(\cref{prop:rg-formation-endogenous-cycle}). Under both laws \(W\) and the
clock have the same law and are independent of the initial state, and given
the past the choices at a ring are independent Bernoulli variables with
probabilities \(\beta_{i,i+1}(\mathbf X_{s-};\pi(s))\) and
\(\beta_{i,i+1}(\mathbf X_{s-};1/2)\). Hence
\(\log(d\Pr/d\mathsf Q)\) on \(\mathcal F^{\mathrm{ch}}_t\) is the log density ratio of the
initial laws plus the sum, over rings in \([0,t]\) and receivers, of the log
ratios of the realized choice probabilities. The initial ratio is
\((2\pi)^{-6}e^{-|\mathbf x-\mathsf m_0|^2/2}\mathsf Z e^{U_S(\mathbf x)/\mathrm T}\)
with \(\mathsf Z=\frac1{16}\sum_S\int e^{-U_S/\mathrm T}d\mathbf x\le(2\pi)^6\),
because \(U_S\ge\mathrm T|\mathbf x|^2/2\); its logarithm is at most
\(U_S(\mathbf x)/\mathrm T\le17|\mathbf x|^2/(2\mathrm T)\), since each of the
four edge Hessians in \(H_S\) has norm \(4\) and \(\|R_{\mathrm{ev}}\|=4\). By
\cref{lem:phil-early-held-seed}(i) each choice ratio is at most
\(\Lambda_{\mathrm L}(t)\le2e^{\epsilon_{\mathrm L}t}\). With \(N_t\) the
number of rings in \([0,t]\), the positive part of \(\log(d\Pr/d\mathsf Q)\) is
at most \(17|\mathbf X_0|^2/(2\mathrm T)+4N_t(\log2+\epsilon_{\mathrm L}t)\),
whose \(L^2(\Pr)\) norm is at most \(\mathfrak L_t\) because
\(\E|\mathbf X_0|^4=(|\mathsf m_0|^2+12)^2+4|\mathsf m_0|^2+24
\le(|\mathsf m_0|^2+14)^2\) and \(\E N_t^2=t+t^2\). The chain rule
\cref{eq:phil-conditioned-parent-chain} with the hidden conditional
\(\nu_{\mathrm T,Z}\) of \(\mu_{\mathrm T,1/2}\), and marginalization to time
\(t\), give
\[
 \mathcal H_{\mathrm{opt}}(q^{\mathscr G_t}_t)
 \le\KL(q^{\mathscr G_t}_t\Vert\mu_{\mathrm T,1/2})
 \le\log\frac1{\Pr(\mathscr G_t)}+
   \frac{\E[\mathbf 1_{\mathscr G_t}\log(d\Pr/d\mathsf Q)]}{\Pr(\mathscr G_t)}
 \le\log\frac1{\Pr(\mathscr G_t)}+\frac{\mathfrak L_t}{\Pr(\mathscr G_t)^{1/2}},
\]
the last step by the Cauchy--Schwarz inequality. Since \(x\log(1/x)\)
increases on \((0,1/e]\), the display of (d) follows. The physical bound uses
\cref{lem:phil-early-physical-comparison}; the native configuration has
finite second moments, and conditioning multiplies them by at most
\(1/\Pr(\mathscr E^\circ_\theta)\).

For the limit, let
\(t\le\epsilon^{-4}/(663552\log(1/\epsilon))\). Then
\(\epsilon^{-4}/(36864t)\ge18\log(1/\epsilon)\), so
\(\mathsf n_t\le24\epsilon^{18}\), and
\(4\Lambda_{\mathrm L}(t)te^{-1/\mathrm T}\le8te^te^{-1/\mathrm T}\) is
smaller than any power of \(\epsilon\). Also
\(\mathfrak L_t=O(\epsilon^{-8})\) for fixed \(\mathsf m_0\), because
\(t\le\epsilon^{-4}\). Hence \((p^\circ_t)^{1/2}\mathfrak L_t=O(\epsilon)\) and
\(p^\circ_t\log(1/p^\circ_t)\to0\), while \(\Pr(\mathscr B^\circ)>0\) does not
depend on \(\epsilon\).

(e) On \(\{S_s=M\text{ for }s\le\theta\}\) one has
\(\mathbf X_s-e^{-K_Ms}\mathbf X_0=\mathbf Z^M_s\), which gives the inclusion
and shows that each piece of \(\mathscr E^{\mathrm{path}}_\theta\) is the
intersection of
\(\mathscr B^{\mathrm{path}}=\{|\mathbf X_0-\sigma\mathbf x^\circ|<\varrho_\circ\}
\cap\{S_s=M\text{ for }s\le\frac14\}\cap\{\sup_{s\le1/4}|\mathbf Z^M_s|\le\delta_\epsilon\}\),
which lies in \(\mathcal F_{1/4}\), with
\(\mathscr A^{\mathrm{path}}=\{S_s=M\text{ for }\frac14<s\le\theta\}\cap
\{\sup_{s\le\theta}|\mathbf Z^M_s|\le\delta_\epsilon\}\). On
\(\mathscr B^{\mathrm{path}}\), \cref{lem:phil-early-held-seed}(iii) gives
\(|F_M^\top\mathbf X_{1/4}|\le\varrho_\circ+\delta_\epsilon\), and the Brownian
increments after \(1/4\) are independent of \(\mathcal F_{1/4}\), so part
(iv) of that lemma with \(s=1/4\) bounds the discarded information of the held
continuation at \(\theta\) by \(\eta^{\mathrm{path}}(\theta)\); on
\(\mathscr B^{\mathrm{path}}\cap\mathscr A^{\mathrm{path}}\) the native
process equals that continuation. Let \(\mathscr O\) be the event
\(\{|\mathbf X_0-\sigma\mathbf x^\circ|<\varrho_\circ\}\cap\{S_0=M\}\). Since
\(W\) is independent of the initial state and of the clock,
\(\Pr(\mathscr B^{\mathrm{path}})\ge\Pr(\mathscr O)e^{-1/4}(1-\mathsf n_\theta)\),
the clock-based seed intersected with the tube on \([0,1/4]\) lying in
\(\mathscr B^{\mathrm{path}}\), while
\(\Pr(\mathscr B^{\mathrm{path}}\cap\{\text{tube fails on }[0,\theta]\})
\le\Pr(\mathscr O)\mathsf n_\theta\). Departures in \((1/4,\theta]\) are bounded
as in (a), after the \(\mathcal F_{1/4}\)-measurable seed. Hence
\(\Pr((\mathscr A^{\mathrm{path}})^c\mid\mathscr B^{\mathrm{path}})
\le p^{\mathrm{path}}_\theta\), and
\cref{lem:phil-early-discarded-information}(i), convexity over the four
pieces, and \cref{lem:phil-early-physical-comparison} finish the proof.
\(\square\)

\propositionheading{Activation at the first qualifying completion}{prop:phil-early-first-completion}
In \cref{def:phil-early-seeded-prefix}, let \(\epsilon_{\mathrm L}\in[0,1]\),
and let \(E_{\mathrm T}\) and \(K^{\mathrm{blk}}_{\mathrm T}\) be the contrast event and the
first-block mark of \cref{prop:phil-early-block-null} with its regime
parameter equal to \(\mathrm T^{1/8}=\epsilon\). Let \(\mathsf t_{\mathrm c}\)
be the least, over eligible blocks, of the infima of their qualifying
completion times, so that \(\mathsf t_{\mathrm c}<\infty\) exactly on
\(E_{\mathrm T}\). With \(p^\circ\) and \(\mathrm c_{\mathsf m}\) from
\cref{prop:phil-early-seeded-parent}, put
\[
 \mathscr U^\circ=\bigcup_{\sigma,M}\bigl(\mathscr B_{\sigma,M}\cap
 \mathscr A^M_{\mathsf t_{\mathrm c}\wedge\theta_\circ}\bigr),
 \qquad
 p_*=p^\circ_{t^\circ_*},
 \qquad
 \mathrm a_\circ=\min\{\epsilon,\varrho_\circ/4\}.
\]

(a) \emph{Completion on the seeded prefix.} For every
\(\theta\in[\theta_\circ,1/(12\mathrm T)]\),
\(\mathscr E^\circ_\theta\subseteq\mathscr U^\circ\subseteq E_{\mathrm T}\).
On \(\mathscr U^\circ\cap\mathscr B_{\sigma,M}\), the mark \(K^{\mathrm{blk}}_{\mathrm T}\)
is one of the two matched pairs of \(M\),
\(\mathsf t_{\mathrm c}\le\theta_\circ\), the assignment equals \(M\) and the
native path equals \(\widetilde{\mathbf X}^M\) on
\([0,\mathsf t_{\mathrm c}]\), and if moreover
\(|n^-_e\cdot\mathbf X_0|\le\mathrm a_\circ\) for an edge \(e\) of \(M\), then
\(\mathsf t_{\mathrm c}\le\frac14\log(1/\epsilon)+\frac{23}{20}\).
Moreover \(\Pr(\mathscr U^\circ)\ge\mathsf c(\mathsf m_0)\).

(b) \emph{Divergence at completion.} Let \(\mathscr G\) be
\(\mathscr U^\circ\) or \(\mathscr E^\circ_\theta\) with
\(\theta\in[\theta_\circ,t^\circ_*]\), and let
\(q^{\mathscr G}=\operatorname{Law}(\mathbf X_{\mathsf t_{\mathrm c}},
S_{\mathsf t_{\mathrm c}}\mid\mathscr G)\). Then
\[
 \mathcal H_{\mathrm{opt}}(q^{\mathscr G})\ge
 e^{-10\mathrm T\theta_\circ}\max\Bigl\{
 \frac{e^{-8}\varrho_\circ^2}{14}(\mathrm c_{\mathsf m}-7p_*)\mathrm T^{-1/2},
 \frac{e^{-46/5}}2\bigl(0.1427\mathrm c_{\mathsf m}\varrho_\circ\mathrm a_\circ
   -\varrho_\circ^2p_*\bigr)\epsilon^{-6}\Bigr\}
 -65\theta_\circ-\log2,
\]
and \(\mathcal H_{\mathrm{phys}}(q^{\mathscr G})\ge
\mathcal H_{\mathrm{opt}}(q^{\mathscr G})-1/2000\). For each fixed
\(\mathsf m_0\), both are at least
\(7.2\cdot10^{-8}\mathrm c_{\mathsf m}\mathrm T^{-5/8}(1+o(1))\) as
\(\epsilon\to0\).

(c) \emph{The whole contrast event.}
\[
 \mathcal H_{\mathrm{opt}}\bigl(\operatorname{Law}(\mathbf X_{\mathsf t_{\mathrm c}},
 S_{\mathsf t_{\mathrm c}}\mid E_{\mathrm T})\bigr)
 \ge\mathsf c(\mathsf m_0)\mathcal H_{\mathrm{opt}}(q^{\mathscr U^\circ})-\log2,
\]
the same holds for \(\mathcal H_{\mathrm{phys}}\), and both tend to infinity
as \(\epsilon\to0\) for each fixed \(\mathsf m_0\).
\status{ESTABLISHED}

This is a refutation for the completion clock: activated at
\(\mathsf t_{\mathrm c}\) from the recognition law conditioned on the seeded
prefix, or on the whole contrast event, no parent of the type of
\cref{prop:rg-formation-endogenous-cycle}(b) has a vanishing conditional error,
in either arm and at every mobility in \([0,1]\). The weight
\(\mathsf c(\mathsf m_0)\) is tiny, so the bound of part (c) becomes positive only
at very small \(\epsilon\); it is a statement about the limit. The first
completion occurs by \(\theta_\circ\), well before the threshold
\(\log(1/\epsilon)\) of \cref{prop:phil-early-seeded-parent}. The order
\(\mathrm T^{-5/8}\) comes from seeds whose first pair forms near
\(\frac14\log(1/\epsilon)\) while the other matched pair, whose hidden
discrepancy the channel also discards, still carries an imprint of order
\(\varrho_\circ\epsilon\). At \(\mathsf t_{\mathrm c}\) only lower bounds are
proved.

\noindent\emph{Proof.}
(a) The time \(\mathsf t_{\mathrm c}\) is a measurable function of the path,
by the measurability argument in the proof of
\cref{prop:phil-early-block-null}, \(\mathbf X\) is continuous, and \(S\) is
right-continuous and piecewise constant, so
\((\mathbf X_{\mathsf t_{\mathrm c}},S_{\mathsf t_{\mathrm c}})\) is measurable
on \(E_{\mathrm T}\). Fix a seed and let \(s\ge1\) satisfy
\(\mathrm Ts\le1/12\). On \(\mathscr B_{\sigma,M}\cap\mathscr A^M_s\) the
native path is held on \([0,s]\), and every eligible block other than the two
matched pairs of \(M\), namely the other two pairs and the four triples,
contains an edge exterior to \(M\). For a completion time \(t\in[1,s+1)\) the
window \([t-1,t]\) meets \([1/4,s]\), where that edge's belief relative
entropy exceeds \(1/2>\epsilon^4\) by \cref{lem:phil-early-held-seed}(iii), so
such a block has no qualifying completion time in \([1,s+1)\). If
\(s\ge\theta_\circ\), the same lemma shows that both matched pairs are
eligible, have internal relative entropies at most \(\epsilon^4\) on
\([\theta_\circ-1,\theta_\circ]\), since
\(\theta_\circ-1=\frac12\log(1/\epsilon)\ge1/4\), and have boundary belief
relative entropies above \(1/2\) there; as \(\theta_\circ\le2/\epsilon\), they
qualify at \(\theta_\circ\).

These facts give the assertions. The inclusion
\(\mathscr E^\circ_\theta\subseteq\mathscr U^\circ\) holds because
\(\mathscr A^M_\theta\subseteq\mathscr A^M_{\mathsf t_{\mathrm c}\wedge\theta_\circ}\).
On \(\mathscr U^\circ\cap\mathscr B_{\sigma,M}\), if
\(\mathsf t_{\mathrm c}>\theta_\circ\), then \(\mathscr A^M_{\theta_\circ}\)
holds and a matched pair qualifies at \(\theta_\circ\), a contradiction; so
\(\mathsf t_{\mathrm c}\le\theta_\circ\), \(E_{\mathrm T}\) holds, the
remainder holds up to \(\mathsf t_{\mathrm c}\), and the path is held there by
the pathwise identification in the proof of
\cref{prop:phil-early-seeded-parent}(a). Applying the first paragraph with
\(s=\mathsf t_{\mathrm c}\), which is at least \(1\) because a completion time
ends a duration-one interval, every block containing an exterior edge has
infimum at least \(\mathsf t_{\mathrm c}+1\), so the least infimum is attained
by a matched pair of \(M\), and \(K^{\mathrm{blk}}_{\mathrm T}\) is one of them. On the
subevent \(|n^-_e\cdot\mathbf X_0|\le\mathrm a_\circ\), write
\(\mathbf d_0=\mathbf X_0-\sigma\mathbf x^\circ\) for the offset. Then
\(\sqrt2n^+_e\cdot\mathbf X_0=-\sigma+\sqrt2n^+_e\cdot\mathbf d_0\) and
\(n^-_e\cdot\mathbf X_0=n^-_e\cdot\mathbf d_0\) by
\cref{lem:phil-early-held-seed}(ii), these coordinates decay at the rates
\(8+\mathrm T\) and \(4+5\mathrm T\) along the held flow, and both residuals
of \(e\) satisfy
\(|r(\widetilde{\mathbf X}^M_u)|\le(1+\sqrt2\varrho_\circ)e^{-8u}
+\mathrm a_\circ e^{-4u}+\epsilon^2/8\). For
\(u\ge\frac14\log(1/\epsilon)+\frac3{20}\) this is at most
\([(1+\sqrt2\varrho_\circ)e^{-6/5}+e^{-3/5}+\frac18]\epsilon^2<0.98\epsilon^2\),
so the pair of \(e\) has internal relative entropies at most \(\epsilon^4\)
from then on. Since \(\frac14\log(1/\epsilon)+\frac{23}{20}\le\theta_\circ\),
the same contradiction argument gives
\(\mathsf t_{\mathrm c}\le\frac14\log(1/\epsilon)+\frac{23}{20}\). Finally
\(\Pr(\mathscr U^\circ)\ge\Pr(\mathscr E^\circ_{\theta_\circ})
\ge(1-p^\circ_{\theta_\circ})2\mathsf c(\mathsf m_0)\) by
\cref{prop:phil-early-seeded-parent}(a), and \(p^\circ_{\theta_\circ}\le
p_*<1/100\), because \(73728\epsilon^4t^\circ_*<0.008\) as in the proof of
\cref{prop:rg-formation-endogenous-cycle} and
\(4\Lambda_{\mathrm L}(t^\circ_*)t^\circ_*e^{-1/\mathrm T}
\le8e\epsilon^{-2}t^\circ_*e^{-1/\mathrm T}<10^{-3}\).

(b) Fix a seed with edges \(e,f\) of \(M\), write
\(\lambda_2=4+5\mathrm T\), and let \(\xi_k(u)\) be the hidden coordinates of
\(\mathbf Z^M_u\). Under \(\Pr(\cdot\mid\mathscr B_{\sigma,M})\) they are
scalar Ornstein--Uhlenbeck noises driven by independent Brownian motions, and
the maximal estimate of the paragraph \emph{Episode} of the proof of
\cref{prop:rg-formation-endogenous-cycle} gives
\(\E\sup_{u\le\theta_\circ}\xi_k(u)^2\le32\mathrm T\theta_\circ\). On
\(\mathscr G\cap\mathscr B_{\sigma,M}\), part (a) gives
\(z_k(\mathbf X_{\mathsf t_{\mathrm c}})=e^{-\lambda_k\mathsf t_{\mathrm c}}
z_k(\mathbf X_0)+\xi_k(\mathsf t_{\mathrm c})\) with
\(\mathsf t_{\mathrm c}\le\theta_\circ\); also
\(\Pr(\mathscr G^c\mid\mathscr B_{\sigma,M})\le p_*\), by
\cref{prop:phil-early-seeded-parent}(a) and the monotonicity of \(p^\circ\) in
\(\theta\). With \((x+y)^2\ge x^2/2-y^2\),
\(|z_k(\mathbf X_0)|<\varrho_\circ\), and the moment bounds of the proof of
\cref{prop:phil-early-seeded-parent}(c),
\[
 \E\bigl[|z_\perp(\mathbf X_{\mathsf t_{\mathrm c}})|^2\mathbf 1_{\mathscr G}
   \bigm|\mathscr B_{\sigma,M}\bigr]
 \ge\tfrac12e^{-2\lambda_2\theta_\circ}\varrho_\circ^2
   \bigl(\mathrm c_{\mathsf m}/7-p_*\bigr)-64\mathrm T\theta_\circ,
\]
and \(e^{-2\lambda_2\theta_\circ}=\epsilon^4e^{-8}e^{-10\mathrm T\theta_\circ}\).
For the second bound keep only \(z_3\) on the subevent
\(\{|z_2(\mathbf X_0)|\le\mathrm a_\circ\}\), where part (a) gives
\(\mathsf t_{\mathrm c}\le\theta_{\mathrm r}=\frac14\log(1/\epsilon)+\frac{23}{20}\)
and \(e^{-2\lambda_2\theta_{\mathrm r}}\ge
\epsilon^2e^{-46/5}e^{-10\mathrm T\theta_\circ}\):
\[
 \E\bigl[z_3(\mathbf X_{\mathsf t_{\mathrm c}})^2\mathbf 1_{\mathscr G}
   \mathbf 1_{\{|z_2(\mathbf X_0)|\le\mathrm a_\circ\}}\bigm|\mathscr B_{\sigma,M}\bigr]
 \ge\tfrac12e^{-2\lambda_2\theta_{\mathrm r}}
   \bigl(\E[z_3(\mathbf X_0)^2\mathbf 1_{\{|z_2(\mathbf X_0)|\le\mathrm a_\circ\}}
     \mid\mathscr B_{\sigma,M}]-\varrho_\circ^2p_*\bigr)
   -32\mathrm T\theta_\circ .
\]
For the uniform law on the twelve-dimensional ball of radius
\(\varrho_\circ\), one orthonormal coordinate divided by \(\varrho_\circ\) has
density \(\mathrm c_{12}(1-x^2)^{11/2}\) on \([-1,1]\), with
\(\mathrm c_{12}=\Gamma(7)/(\sqrt\pi\Gamma(13/2))=46080/(10395\pi)>1.411\),
and given that coordinate the other eleven are uniform on a ball of radius
\(\varrho_\circ(1-x^2)^{1/2}\), so a second orthonormal coordinate has
conditional second moment \(\varrho_\circ^2(1-x^2)/13\). Write
\(\E_{\mathrm{unif}}\) for the expectation when the offset
\(\mathbf d_0=\mathbf X_0-\sigma\mathbf x^\circ\) of (a) has this uniform law.
Since \(\mathrm a_\circ/\varrho_\circ\le1/4\),
\[
 \E_{\mathrm{unif}}\bigl[z_3(\mathbf d_0)^2
   \mathbf 1_{\{|z_2(\mathbf d_0)|\le\mathrm a_\circ\}}\bigr]
 =\frac{\mathrm c_{12}\varrho_\circ^2}{13}
   \int_{-\mathrm a_\circ/\varrho_\circ}^{\mathrm a_\circ/\varrho_\circ}
   (1-x^2)^{13/2}dx
 \ge\frac{2\mathrm c_{12}}{13}\Bigl(\frac{15}{16}\Bigr)^{13/2}
   \varrho_\circ\mathrm a_\circ
 \ge0.1427\varrho_\circ\mathrm a_\circ,
\]
and the density bound \(\mathrm c_{\mathsf m}\) of the proof of
\cref{prop:phil-early-seeded-parent}(c) transfers it to the seed. The law
\(q^{\mathscr G}\) is a mixture of the four seed laws and is carried by the
matching strata, so the Donsker--Varadhan bound of that proof, with the factor
\((4+5\mathrm T)/(4\mathrm T)\ge1/\mathrm T\) applied when the bracket is
positive and \(64\theta_\circ(1+5\mathrm T/4)\le65\theta_\circ\), gives the
display; when both brackets are negative the display is negative and holds
trivially. On \(\mathscr G\) the hidden coordinates have norm at most
\(\varrho_\circ+\delta_\epsilon\), and
\cref{lem:phil-early-physical-comparison} gives the physical bound. For the
limit, \(\mathrm a_\circ=\epsilon\) once \(\epsilon\le1/400\),
\(p_*/\epsilon\to0\), and \(\theta_\circ=O(\log(1/\epsilon))\), so the second
bracket gives
\(\frac12e^{-46/5}0.1427\varrho_\circ\mathrm c_{\mathsf m}\epsilon^{-5}(1+o(1))\)
with \(\frac12e^{-46/5}0.1427\varrho_\circ>7.2\cdot10^{-8}\) and
\(\epsilon^{-5}=\mathrm T^{-5/8}\).

(c) By (a), \(\mathscr U^\circ\subseteq E_{\mathrm T}\) and
\(\Pr(\mathscr U^\circ\mid E_{\mathrm T})\ge\Pr(\mathscr U^\circ)\ge
\mathsf c(\mathsf m_0)\). The law given \(E_{\mathrm T}\) is the mixture of
the laws given \(\mathscr U^\circ\) and given
\(E_{\mathrm T}\setminus\mathscr U^\circ\), and so are its evaluator lifts,
so \cref{lem:phil-early-discarded-information}(i) gives both inequalities.
The limits follow from (b). \(\square\)

\propositionheading{Parent dynamics after the seeded early prefix}{prop:phil-early-parent-dynamics}
In \cref{def:phil-early-seeded-prefix}, let \(\epsilon_{\mathrm L}\in[0,1]\),
let \(\theta\ge1/4\) with \(\mathrm T\theta\le1/12\) and
\(p^\circ_\theta<1/2\) as in \cref{prop:phil-early-seeded-parent}, and let
\(\mathsf h>0\) with \(\mathrm T(\theta+\mathsf h)\le1/10\). Activate the parent
of \cref{prop:phil-actual-parent-dynamics}, with the kernel
\cref{eq:phil-actual-parent-kernel} and own-law learning
\cref{eq:phil-actual-parent-learning} at mobility \(\epsilon_{\mathrm L}\), from
\(q^C_0=\mu_{\mathrm{act}}=C_\#q^\circ_\theta\) and \(p^C(0)=p(\theta)\).

(i) \(\mu_{\mathrm{act}}\) is carried by the union over seeds of
\(\{M\}\times\{y\in W_M:|y-P_{W_M}\bar{\mathbf x}^{\sigma,M}(\theta)|
\le\varrho_\circ+\delta_\epsilon\}\).

(ii) The coupled parent system has exactly one solution on \([0,\infty)\),
and its rows satisfy
\(p^C_i(u)\in[e^{-\epsilon_{\mathrm L}(\theta+u)}/2,
1-e^{-\epsilon_{\mathrm L}(\theta+u)}/2]\).

(iii) Let \(\mathsf P^{\mathrm{fine}}_{\theta,\mathsf h}\) be the law of
\((C(\mathbf X_{\theta+u},S_{\theta+u}))_{0\le u\le\mathsf h}\) conditioned on
\(\mathscr E^\circ_\theta\), the native rows continuing on their own path, and
\(\mathsf P^C_{\mathsf h}\) the parent path law. Then
\[
 \operatorname{TV}\bigl(\mathsf P^{\mathrm{fine}}_{\theta,\mathsf h},
                        \mathsf P^C_{\mathsf h}\bigr)
 \le24\exp\Bigl[-\frac1{76800\mathrm T\mathsf h}\Bigr]
 +4\bigl(1+\sqrt{8/3}\bigr)\Lambda_{\mathrm L}(\theta+\mathsf h)
   \mathsf he^{-1/\mathrm T}.
\]
The comparison includes every retained coordinate and the assignment path,
hence the loop. The right side does not depend on \(\mathsf m_0\) or on
\(\Pr(\mathscr E^\circ_\theta)\), and it tends to zero as \(\epsilon\to0\),
uniformly in \(\theta\), \(\epsilon_{\mathrm L}\in[0,1]\), and
\(\mathsf m_0\), whenever \(\mathrm T\mathsf h\to0\), for instance for
\(\mathsf h=\mathrm T^{-1+\mathrm q}\) with \(\mathrm q\in(0,1)\).

In particular, if \(\epsilon\to0\) with \(\theta-\log(1/\epsilon)\to\infty\),
\(\epsilon^4\theta\to0\), and \(\mathrm T\mathsf h\to0\), for instance for
\(\theta\in[t_\epsilon,t^\circ_*]\), then
\(\mathcal H_{\mathrm{opt}}(q^\circ_\theta)\),
\(\mathcal H_{\mathrm{phys}}(q^\circ_\theta)\), and the total variation of
(iii) tend to zero separately, uniformly in \(\mathsf m_0\) and
\(\epsilon_{\mathrm L}\in[0,1]\).
\status{ESTABLISHED}

\noindent\emph{Proof.}
(i) follows from \cref{prop:phil-early-seeded-parent}(a), since \(P_{W_M}\)
is a contraction.

(ii) The argument of the proof of \cref{prop:phil-actual-parent-dynamics},
by coupled prescribed-row constructions, contraction of the
variation-of-constants map on short intervals, and the lower bounds
\(p^C_i(u)\ge p^C_i(0)e^{-\epsilon_{\mathrm L}u}\) and
\(1-p^C_i(u)\ge(1-p^C_i(0))e^{-\epsilon_{\mathrm L}u}\), uses only an
activation law with finite second moment and interior initial rows; neither
the selection tolerances nor \(d<1/4\) enter it. Here \(\mu_{\mathrm{act}}\)
has bounded support by (i), and \(p(\theta)\) has entries in
\([e^{-\epsilon_{\mathrm L}\theta}/2,1-e^{-\epsilon_{\mathrm L}\theta}/2]\) by
\cref{lem:phil-early-held-seed}(i); the two lower bounds give the stated
interval. For \(\epsilon_{\mathrm L}=0\) the rows stay at \(p(\theta)\).

(iii) \emph{Coupling.} Since the native rows are deterministic, the native
path after \(\theta\) is a measurable function of
\((\mathbf X_\theta,S_\theta)\) and of the inputs after \(\theta\), which are
independent of \(\mathcal F_\theta\ni\mathscr E^\circ_\theta\). On one
probability space take independent objects: a state
\((\mathbf X^N_0,S^N_0)\) with law \(q^\circ_\theta\), a standard Brownian
motion \(W'\), a rate-one clock, uniform variables for each ring and receiver,
and draws \(\zeta_k\sim\nu_{\mathrm T,Z}\), one for each ring. The native copy
starts at \((\mathbf X^N_0,S^N_0)\), solves
\(d\mathbf X^N=-K_{S^N}\mathbf X^Ndu+\sqrt{2\mathrm T}dW'\) between rings, and
redraws with rows \(p(\theta+u)\) at \(\mathbf X^N_{u-}\); its retained path
has law \(\mathsf P^{\mathrm{fine}}_{\theta,\mathsf h}\). The parent copy
starts at \(C(\mathbf X^N_0,S^N_0)\), solves
\(dY^C=-K_SY^Cdu+\sqrt{2\mathrm T}P_{W_S}dW'\) on stratum \(S\), and at the
\(k\)th ring redraws with rows \(p^C(u)\) at
\(\widetilde{\mathbf x}=Y^C_{u-}+F_S\zeta_k\) on a matching stratum
(\(\widetilde{\mathbf x}=Y^C_{u-}\) otherwise), moving to
\((S',P_{W_{S'}}\widetilde{\mathbf x})\); since \(P_{W_S}W'\) is a standard
Brownian motion on \(W_S\) and \(p^C\) is deterministic, its path law is
\(\mathsf P^C_{\mathsf h}\). Let \(M_0=S^N_0\) and let \(\sigma_{\mathrm D}\)
be the first ring at which either copy proposes an assignment other than
\(M_0\). Before \(\sigma_{\mathrm D}\), \(K_{M_0}\) commutes with
\(P_{W_{M_0}}\), so \(P_{W_{M_0}}\mathbf X^N\) and \(Y^C\) solve the same
linear equation with the same noise and the same initial value, and a ring
that keeps \(M_0\) moves neither copy, since
\(P_{W_{M_0}}(Y^C+F_{M_0}\zeta_k)=Y^C\). The retained paths therefore agree on
\([0,\mathsf h]\) when \(\sigma_{\mathrm D}>\mathsf h\), and the total
variation is at most \(\Pr(\sigma_{\mathrm D}\le\mathsf h)\).

\emph{Tube and gap.} Let
\(\mathbf Z'_u=\sqrt{2\mathrm T}\int_0^ue^{-K_{M_0}(u-v)}dW'_v\) and
\(\mathsf t_\Gamma=\inf\{u:|\mathbf Z'_u|\ge1/20\}\). Since \(W'\) is independent
of the initial state, the argument of \cref{prop:phil-early-seeded-parent}(a)
with radius \(1/20\) in place of \(\delta_\epsilon\), and with the reflection
principle for the closed event, gives
\(\Pr(\mathsf t_\Gamma\le\mathsf h)\le24\exp[-(1/20)^2/(192\mathrm T\mathsf h)]
=24\exp[-1/(76800\mathrm T\mathsf h)]\). By
\cref{prop:phil-early-seeded-parent}(a),
\(|\mathbf X^N_0-\bar{\mathbf x}(\theta)|\le\varrho_\circ+\delta_\epsilon\)
for some sign \(\sigma\), where
\(\bar{\mathbf x}=\bar{\mathbf x}^{\sigma,M_0}\). For
\(u\le\sigma_{\mathrm D}\wedge\mathsf t_\Gamma\) the native configuration before the
ring is \(e^{-K_{M_0}u}\mathbf X^N_0+\mathbf Z'_u\), and
\(\bar{\mathbf x}(\theta+u)=e^{-K_{M_0}u}\bar{\mathbf x}(\theta)\), so
\(|\mathbf X^N_u-\bar{\mathbf x}(\theta+u)|\le\mathrm R=
\varrho_\circ+\delta_\epsilon+1/20<0.0601\). For an exterior edge \(g\), the
common residual \(c_g=(r^b_g+r^m_g)/2\) has a covector of norm \(\sqrt2\) that
annihilates \(Z_{M_0}\) (proof of \cref{prop:rg-formation-endogenous-cycle}),
\(\mathcal D_g\ge c_g^2\), and along \(\bar{\mathbf x}(t)\),
\(|c_g|=e^{-\mathrm Tt}(3/2-e^{-8t}/2)\) by
\cref{lem:phil-early-held-seed}(ii). For \(t\in[\theta,\theta+\mathsf h]\),
with \(\theta\ge1/4\) and \(\mathrm T(\theta+\mathsf h)\le1/10\), this gives
\(|c_g(\mathbf X^N_u)|\ge e^{-1/10}(3/2-e^{-2}/2)-\sqrt2\mathrm R>1.21\). For a
selected edge, \(\mathcal D_e=(\sqrt2n^+_e\cdot\mathbf x)^2+(n^-_e\cdot\mathbf x)^2\)
is the square of a seminorm of operator norm \(\sqrt2\) that equals
\(e^{-(8+\mathrm T)t}\le e^{-2}\) at \(\bar{\mathbf x}(t)\), so
\(\mathcal D_e(\mathbf X^N_u)\le(e^{-2}+\sqrt2\mathrm R)^2\le2(e^{-2}+\mathrm R)^2<0.08\).
Likewise \(z_+\) is \(1\)-Lipschitz, retained, and equal to
\(-\sigma e^{-(8+\mathrm T)t}\) on the center path, so
\(\frac43z_+^2\le2(e^{-2}+\mathrm R)^2<0.08\). Since \(1.21^2-0.08>1\),
the energy gap exceeds one for both copies.

\emph{Hazards.} At a ring at relative time
\(u\le\sigma_{\mathrm D}\wedge\mathsf t_\Gamma\), the native copy departs with
conditional probability at most
\(\sum_i\beta_{i,\mathrm{ext}}\le4\Lambda_{\mathrm L}(\theta+u)e^{-1/\mathrm T}\)
by \cref{lem:phil-early-held-seed}(i). For the parent,
\(\mathcal D_g(\widetilde{\mathbf x})\ge c_g(Y^C)^2\) and
\(\mathcal D_e(\widetilde{\mathbf x})=(z_++z_1)^2+z_2^2\) with independent
\(z_1\sim\mathcal N(0,\mathrm T/(8+\mathrm T))\) and
\(z_2\sim\mathcal N(0,\mathrm T/(4+5\mathrm T))\), and similarly for \(f\) with
\(-z_1\) and \(z_3\). The identity
\(\E e^{(\mathrm c_1+Z)^2/\mathrm T}=(1-2s^2/\mathrm T)^{-1/2}
\exp[\mathrm c_1^2/(\mathrm T(1-2s^2/\mathrm T))]\), valid for
\(Z\sim\mathcal N(0,s^2)\) with \(s^2<\mathrm T/2\), gives the factors
\(((8+\mathrm T)/(6+\mathrm T))^{1/2}\le\sqrt{4/3}\) and
\(((4+5\mathrm T)/(2+5\mathrm T))^{1/2}\le\sqrt2\) and an exponent at most
\(\frac43z_+^2/\mathrm T\). With the row bound of (ii), the parent departs with
conditional probability at most
\(4\sqrt{8/3}\Lambda_{\mathrm L}(\theta+u)e^{-1/\mathrm T}\). The variables of
a ring are independent of its past, so the number \(N_{\mathrm D}\) of
departing rings in \([0,\mathsf h\wedge\sigma_{\mathrm D}\wedge\mathsf t_\Gamma]\)
has, by the compensator of the clock,
\(\E N_{\mathrm D}\le4(1+\sqrt{8/3})\Lambda_{\mathrm L}(\theta+\mathsf h)
\mathsf he^{-1/\mathrm T}\). If \(\sigma_{\mathrm D}\le\mathsf h\), then
\(\mathsf t_\Gamma\le\mathsf h\) or \(N_{\mathrm D}\ge1\), which proves the bound.

\emph{Limits.} The first term depends only on \(\mathrm T\mathsf h\). For
\(\epsilon_{\mathrm L}\le1\),
\(\Lambda_{\mathrm L}(\theta+\mathsf h)\le2e^{1/(10\mathrm T)}\) and
\(\mathsf h\le1/(10\mathrm T)\), so the second term is at most
\(3e^{-9/(10\mathrm T)}/\mathrm T\). For \(\mathsf h=\mathrm T^{-1+\mathrm q}\),
\(\mathrm T(\theta+\mathsf h)\le1/10\) once \(\epsilon\) is small. The last
paragraph of the statement combines (iii) with
\cref{prop:phil-early-seeded-parent}(b), whose bound does not depend on
\(\mathsf m_0\), and the window \([t_\epsilon,t^\circ_*]\) satisfies
\(\theta-\log(1/\epsilon)\ge\log(1/\epsilon)\). \(\square\)

\remarkheading{The stationary comparison does not apply to the early prefix}{rem:phil-early-stationary-mismatch}
Let \(\epsilon_{\mathrm L}\in[0,\epsilon]\) and \(0\le t\le t^\circ_*\).
Every row entry \(\pi_{ij}(t)\) lies in \([0.4489,0.5511]\), so every row
assigns more than \(1/4\) to the nonpartner source of either matching and the
hypothesis \(d<1/4\) of \cref{prop:phil-actual-parent-dynamics} fails; and for
each matching \(M\),
\((1-\mu_{\mathrm T,p(t)}(M))/\mu_{\mathrm T,p(t)}(M)\ge0.44\), so the averaged
departure bound \cref{eq:phil-parent-averaged-departure} is at least
\(0.44\) per unit time. Under the hypotheses of
\cref{prop:phil-early-parent-dynamics}, for each matching \(M\),
\[
 \operatorname{TV}\bigl(\mu_{\mathrm{act}},
   \overline\nu^M_{\mathrm T}\otimes\delta_M\bigr)\ge1-3\cdot10^{-18},
\]
where \(\overline\nu^M_{\mathrm T}=(P_{W_M})_\#\mathcal N(0,\mathrm TK_M^{-1})\)
is the retained Gaussian of that proposition. Consequently each leg of the
triangle inequality through the stationary held-matching path law in its proof
exceeds \(1-3\cdot10^{-18}\).
\status{ESTABLISHED}

\noindent\emph{Proof.}
By \cref{lem:phil-early-held-seed}(i),
\(|\pi_{ij}(t)-1/2|\le(1-e^{-\epsilon t^\circ_*})/2\), and
\(\epsilon t^\circ_*<0.1022\) as in the proof of
\cref{prop:phil-seeded-contrast-formation}. The rotation \(i\mapsto i+1\) acts
by an orthogonal map carrying \(U_{S^{(0)}}\) to \(U_{S^{(1)}}\), so the two
matchings have equal Gaussian masses \(\int e^{-U_M/\mathrm T}d\mathbf x\).
Dropping the other fourteen sheets from the normalizer gives
\(\mu_{\mathrm T,p}(M)\le\pi_M/(\pi_M+\pi_{M'})\), where \(M'\) is the other
matching, so
\((1-\mu_{\mathrm T,p}(M))/\mu_{\mathrm T,p}(M)\ge\pi_{M'}/\pi_M
=\prod_i\pi_{i,M'_i}/\pi_{i,M_i}\ge(0.4489/0.5511)^4>0.44\).

For the second claim, \(R_{\mathrm{ev}}=2\mathfrak D^\top\mathfrak D\) has
eigenvalue \(4\) with multiplicity four and \(0\) with multiplicity eight,
because \(\mathfrak D\mathfrak D^\top=2I_4\). It leaves \(\ker H_M\) and the
range of \(H_M\) invariant, and on the range it has eigenvalues \(0,0,4,4\)
(\(R_{\mathrm{ev}}n^+_g=0\), \(R_{\mathrm{ev}}n^-_g=4n^-_g\)), so on
\(\ker H_M\subseteq W_M\) it has eigenvalue \(4\) twice and \(0\) six times.
Under \(\mathcal N(0,\mathrm TK_M^{-1})\) the component in \(\ker H_M\) has
covariance \((I+R_{\mathrm{ev}})^{-1}\) there, with eigenvalues \(1\) six times
and \(1/5\) twice, hence a density at most \(5(2\pi)^{-4}\) on this
eight-dimensional space. The projection onto \(\ker H_M\) of a retained ball of
radius \(\mathrm R_0=\varrho_\circ+\delta_\epsilon<0.010008\) lies in an
eight-dimensional ball of volume \(\pi^4\mathrm R_0^8/24\). By (i) of
\cref{prop:phil-early-parent-dynamics}, \(\mu_{\mathrm{act}}\) is carried by
a set whose part on the stratum of \(M\) consists of two such balls, of
\(\overline\nu^M_{\mathrm T}\)-mass at most
\(2\cdot5(2\pi)^{-4}\pi^4\mathrm R_0^8/24=10\mathrm R_0^8/384<3\cdot10^{-18}\),
and the rest of that set has \(\overline\nu^M_{\mathrm T}\otimes\delta_M\)-mass
zero. Both \(\mathsf P^{\mathrm{fine}}_{\theta,\mathsf h}\) and
\(\mathsf P^C_{\mathsf h}\) start from \(\mu_{\mathrm{act}}\), the stationary
path law starts from \(\overline\nu^M_{\mathrm T}\otimes\delta_M\), and the
total variation of path laws is at least that of their initial laws.
\(\square\)

The mismatch is structural. Early departures are rare because the exterior
common residuals stay near \(3/2\) while the partner costs are small, a
pointwise property of the activation law; under the stationary retained law the
slow coordinates are centered at zero with variance of order one, and the
comparison charges the stationary mass of the other matching as potential
flux. The proof of \cref{prop:phil-early-parent-dynamics} therefore uses the
energy gap along the conditioned path instead.

These statements concern the seeded sub-event of the contrast. The conditional
law given the whole contrast event at a deterministic activation time, which
would need control of contrast histories that enter a matching at a random
positive time or start outside the seed balls, is not treated, and neither is
the comparison of learned and frozen first blocks at
\(\epsilon_{\mathrm L}=\mathrm T^{1/8}\) or of two initial means differing in
their loop coordinate; these remain part of
\cref{open:rg-formation-adaptive-family}. The path coupling gives no
conclusion on windows of order \(\mathrm T^{-1}\) or longer, and the parent
path over a following window describes the persistence of a formed block, not a
second formation.
